\documentclass[12pt]{beamer}
% ---------- doc 1 format: 12pt, lmodern serif, pale-blue title band, triangle bullets ----------
\usepackage{lmodern}
\usefonttheme{serif}
\usepackage{amsmath,amssymb}
\usepackage{booktabs}
\usepackage{pgfplots}
\pgfplotsset{compat=1.18}

\definecolor{lectureblue}{RGB}{0,70,150}
\definecolor{paleblue}{RGB}{221,233,248}
\definecolor{examplegreen}{RGB}{0,120,60}

\setbeamertemplate{navigation symbols}{}
\setbeamertemplate{itemize item}[triangle]
\setbeamertemplate{itemize subitem}[triangle]
\setbeamercolor{itemize item}{fg=lectureblue}
\setbeamercolor{itemize subitem}{fg=lectureblue}
\setbeamercolor{frametitle}{bg=paleblue,fg=black}
\setbeamerfont{frametitle}{size=\large,series=\bfseries}
\setbeamertemplate{frametitle}{%
  \nointerlineskip
  \begin{beamercolorbox}[wd=\paperwidth,ht=2.6ex,dp=1.1ex,leftskip=0.45cm]{frametitle}%
    \usebeamerfont{frametitle}\insertframetitle
  \end{beamercolorbox}}
\setbeamercolor{title}{fg=lectureblue}
\setbeamertemplate{footline}{\hfill\scriptsize\insertframenumber\hspace*{0.3cm}\vspace*{0.15cm}}

\newcommand{\keyterm}[1]{\textcolor{examplegreen}{\textbf{#1}}}
\newcommand{\fitO}{}
\newcommand{\exlists}{\setlength{\itemsep}{0pt}}
\newcommand{\ind}[1]{\hspace*{#1}}
\newcommand{\Xbar}{\overline{X}}

\newcommand{\thetahat}{\hat{\theta}}

\title{Section 6.1}
\subtitle{Point Estimation}
\date{}

\begin{document}

% ======================================================================
\begin{frame}
\titlepage
\end{frame}

% ======================================================================
\begin{frame}[t]{Point Estimation}
\begin{itemize}
\item Goal: use sample data to estimate an unknown parameter
  \textcolor{magenta}{$\theta$} (for example $\mu$, $\sigma$ or $p$).
\item[]
\item A \keyterm{point estimate} of $\theta$ is a single number computed from the
  \textbf{observed sample} $x_1,\dots,x_n$. It is a sensible value for $\theta$.
\item A \keyterm{point estimator} \textcolor{magenta}{$\thetahat$} is the statistic used to get it.
  It is computed from the \textbf{random sample} $X_1,\dots,X_n$, so it is a random variable.
\item[]
\item Example: $\textcolor{magenta}{\hat{\mu}}=\Xbar$ is an estimator of $\mu$;
  $\bar{x}=12.4$ is an estimate.
\item Different estimators can give different estimates, so we need criteria to choose one.
\end{itemize}
\end{frame}

% ======================================================================
\begin{frame}[t]{Common Point Estimators}
\begin{center}
\renewcommand{\arraystretch}{1.6}\setlength{\tabcolsep}{14pt}
\begin{tabular}{lll}
\toprule
Parameter & Estimator & Estimate \\
\midrule
mean $\mu$            & $\hat{\mu}=\Xbar$         & $\bar{x}$ \\
variance $\sigma^2$   & $\hat{\sigma}^2=S^2$      & $s^2$ \\
standard deviation $\sigma$ & $\hat{\sigma}=S$    & $s$ \\
proportion $p$        & $\hat{p}=X/n$             & $x/n$ \\
\bottomrule
\end{tabular}
\end{center}
\vspace{0.3cm}
\begin{itemize}
\item Estimator: from the random sample (capital letters). Estimate: from the observed sample (lowercase).
\item Here $X$ is the number of successes in the $n$ trials.
\item Recall: $S^2=\dfrac{1}{n-1}\displaystyle\sum_{i=1}^{n}(X_i-\Xbar)^2$.
\end{itemize}
\end{frame}

% ======================================================================
%% answers: a) xbar = 96/8 = 12 hours
%%          b) s^2 = 28/7 = 4, s = 2 hours
%%          c) phat = 3/8 = 0.375
\begin{frame}[t]{}\fitO\exlists
\textcolor{blue}{\textbf{Example 1:}}\,\ The lifetimes (in hours) of 8 randomly selected batteries are
\begin{center}
12 \quad 15 \quad 9 \quad 14 \quad 10 \quad 13 \quad 11 \quad 12
\end{center}
\begin{itemize}
\item[a)] Find a point estimate of the mean lifetime $\mu$.\\[0.2cm]
\end{itemize}
\end{frame}

\begin{frame}[t]{}\fitO\exlists
\begin{itemize}
\item[b)] Find point estimates of $\sigma^2$ and $\sigma$.\\[3.6cm]
\item[c)] Find a point estimate of the proportion $p$ of all batteries that last more than 12 hours.\\[0.2cm]
\end{itemize}
\end{frame}

% ======================================================================
\begin{frame}[t]{Unbiased Estimators}
\begin{itemize}
\item A point estimator $\thetahat$ is \keyterm{unbiased} if
\[
\textcolor{blue}{E(\thetahat)=\theta}\quad\text{for every value of }\theta.
\]
\item The \keyterm{bias} of $\thetahat$ is \ \textcolor{red}{$E(\thetahat)-\theta$}.
\item[]
\item Unbiased estimators we already know:
  \begin{itemize}
  \item $\Xbar$ for $\mu$, since $E(\Xbar)=\mu$;
  \item $\hat{p}=X/n$ for $p$, since $E(X)=np$;
  \item $S^2$ for $\sigma^2$. This is why $S^2$ divides by $n-1$.
  \end{itemize}
\item $S$ is \textbf{not} unbiased for $\sigma$, but its bias is small.
\item Among unbiased estimators, prefer the one with the smallest variance.
\end{itemize}
\end{frame}

% ======================================================================
%% answers: a) E(th1)=mu, E(th2)=(mu+2mu+mu)/4=mu, E(th3)=2mu/3 -> th1, th2 unbiased; bias of th3 = -mu/3
%%          b) V(th1)=sigma^2/3 ~ 0.333 sigma^2, V(th2)=(1+4+1)sigma^2/16=0.375 sigma^2 -> prefer th1 = Xbar
\begin{frame}[t]{}\fitO\exlists
\textcolor{blue}{\textbf{Example 2:}}\,\ Let $X_1,X_2,X_3$ be a random sample from a distribution
with mean $\mu$ and variance $\sigma^2$. Consider
{\small
\[
\thetahat_1=\frac{X_1+X_2+X_3}{3},\quad
\thetahat_2=\frac{X_1+2X_2+X_3}{4},\quad
\thetahat_3=\frac{X_1+X_2}{3}.
\]}
\begin{itemize}
\item[a)] Which of these are unbiased estimators of $\mu$? Find the bias of any that are not.\\[2.6cm]
\item[b)] Among the unbiased ones, which would you prefer? Why?\\[0.2cm]
\end{itemize}
\end{frame}

\end{document}
